% ===========================================================================
%  sviluppatore/02-ambiente/01-prerequisiti.tex
% ===========================================================================
\section{Che cosa serve avere installato}\label{sez:prerequisiti}

Per la metà locale: Python, e nient'altro di obbligatorio. Per il sito:
Node.js e un Postgres raggiungibile. Le due metà si avviano in modo
indipendente, e per lavorare sul motore il sito non serve --- senza di esso si
perde solo il marketplace.

\begin{center}
\begin{tabular}{@{}L{34mm}L{26mm}L{78mm}@{}}
\toprule
\textbf{Componente} & \textbf{Versione} & \textbf{Note} \\
\midrule
\rigaTenue
Python & \num{3}.10 o più & Il codice usa \cod{match}, \cod{X | None} e
  le \cod{dataclass} con \cod{slots}. \\
pandas & $<$\,\num{3} & Vincolo reale: il codice precede l'API di pandas~3. \\
\rigaTenue
numpy & $<$\,\num{2} & Stesso motivo. In caso di errori nello stack
  numerico, questi due vincoli sono la prima cosa da provare. \\
Node.js & \num{18} o più & Solo per il sito. \\
\rigaTenue
Postgres & \num{14} o più & Solo per il sito. In sviluppo va bene anche
  un'istanza gestita (Neon). \\
\bottomrule
\end{tabular}
\end{center}

\subsection{Installazione delle dipendenze Python}

\begin{codiceInterno}{console}{Dalla radice del progetto}{}
python -m venv .venv
source .venv/bin/activate        # Windows: .venv\Scripts\activate
pip install -r requirements.txt
\end{codiceInterno}

\noindent
Il file \file{requirements.txt} è diviso in blocchi commentati e dichiara
di che cosa è responsabile ogni dipendenza. Tre di essi vanno guardati:

\begin{description}[leftmargin = 1.6em, style = nextline, itemsep = 4pt]
  \item[Stack numerico] \pkg{numpy}, \pkg{pandas}, \pkg{scipy},
    \pkg{scikit-learn}, \pkg{statsmodels}, \pkg{pyarrow}. Senza
    \pkg{pyarrow} non si leggono i \file{parquet}, cioè non si legge niente.
  \item[Machine learning] \pkg{lightgbm}, \pkg{hmmlearn}, \pkg{shap}. Servono
    agli approcci ML e alle modalità di pesatura \cod{regime\_ml} e
    \cod{regime\_style}. Un sistema senza di esse parte, e quelle voci
    semplicemente non funzionano.
  \item[Sorgenti dati] \pkg{yfinance}, \pkg{requests}, \pkg{fredapi},
    \pkg{python-dotenv}.
\end{description}

\begin{nota}{Il modello locale è opzionale}
Lo Studio può generare il codice di un artefatto a partire da una descrizione
in italiano. Il modello gira in locale via Ollama, che non è nelle
dipendenze Python e va installato a parte. Senza, lo Studio funziona
ugualmente: si perde la generazione da testo, restano l'editor e il grafo a
nodi.
\end{nota}

\subsection{Con Docker, invece}

Nella radice ci sono \file{Dockerfile} e \file{docker-compose.yml}. Avviano la
sola dashboard, con i dati montati come volume.

\begin{codiceInterno}{console}{Dalla radice del progetto}{}
UID=$(id -u) GID=$(id -g) docker compose up --build
# poi: http://localhost:5000
\end{codiceInterno}

\begin{attenzione}{Su Linux servono UID e GID}
Senza passarli, i file che il contenitore scrive nei volumi --- i
\file{parquet} scaricati, i risultati dei backtest --- risultano di un altro
proprietario e dall'esterno non si riescono più a modificare. Su macOS e
Windows il problema non si pone.
\end{attenzione}
